\section{From reconstruction to capacity}\label{sec:necessity}

\begin{lemma}\label{lem:cubic}
Under Assumption~\ref{ass:input}, the increasing function
$g(t)=t+Kt^3$, with $K=1/3000$, satisfies
\begin{equation}\label{eq:cubic}
 g\bigl(m(\mu_1\star\cdots\star\mu_d)\bigr)
       \leq\sum_{j=1}^d g(m(\mu_j))
\end{equation}
for every finite collection of conditionally independent experiments
whose laws lie in $\calC$.
\end{lemma}
\begin{proof}
For two laws put $a=m(\mu)$, $b=m(\nu)$, and
$d_0=m(\mu\star\nu)$. Jensen's inequality and
\eqref{eq:saving} imply
\[
 0\leq d_0\leq a+b-cab(a+b)\leq a+b,\qquad c=1/1000.
\]
Use the first upper bound for the linear term and the second for
the cubic term to obtain
\[
 g(d_0)\leq a+b-cab(a+b)+K(a+b)^3
 =g(a)+g(b)+(3K-c)ab(a+b)=g(a)+g(b).
\]
This step does not require $a+b\leq3$. Closure under products
allows iteration. The empty product is uninformative.
\end{proof}

Let $m_v^{(n)}$ be the information about $\sigma_v$ from its
descendants at distance $n$, and let $m_v=\lim_n m_v^{(n)}$.
The limit exists by data processing and lies in $[0,3]$.
Apply Lemma~\ref{lem:cubic} at finite depth and pass to the limit
through the finitely many children of each vertex. With $y_v=g(m_v)$,
we get
\begin{align}
 y_v&\leq\sum_{w\in\ch(v)}g(r_e m_w)\notag\\
 &=\sum_w r_e y_w-K\sum_w r_e(1-r_e^2)m_w^3\notag\\
 &\leq\sum_w r_e y_w
       \left(1-\frac\kappa2(1-r_e^2)y_w^2\right)\notag\\
 &\leq\sum_w\frac{r_e y_w}
                 {\sqrt{1+\kappa(1-r_e^2)y_w^2}}.
                 \label{eq:recursive-labels}
\end{align}
Here $e=(v,w)$, $y_w\leq(1+9K)m_w$, and
$(1+2t)^{-1/2}\geq1-t$ for $t\geq0$. No limiting posterior law
has been assumed to lie in $\calC$.

We now construct a flow directly from these labels. For each vertex
choose coefficients
\[
 0\leq d_e\leq\frac{r_e}{\sqrt{1+\kappa(1-r_e^2)y_w^2}}
 \quad\text{such that}\quad y_v=\sum_w d_e y_w.
\]
If the upper-bound sum is positive, scale all its coefficients by the
ratio of $y_v$ to that sum; otherwise choose them all zero. Put
$\ell_\rho=1$, $\ell_w=\ell_vd_e$, and
$\theta(e)=\ell_w y_w$. This is a conserved flow of mass
$y_\rho$. With $z_v=\ell_v/R_v$, the coefficient bound gives
\[
 \kappa(1-r_e^2)\left(\frac{\theta(e)}{R_w}\right)^2
 \leq z_v^2-z_w^2.
\]
Summing along any ray gives a bound of $z_\rho^2=1$.
Thus $\sqrt\kappa\,\theta$ is admissible and
\begin{equation}\label{eq:strong-upper}
 C(T)\geq\sqrt\kappa\,g(m_\infty)
          \geq\sqrt\kappa\,m_\infty.
\end{equation}
This proves the upper bound in Theorem~\ref{thm:main}. The construction
also covers zero labels, dead ends, and perfectly transmitting edges.
